\subsubsection{Evaluación preliminar del reconocimiento de escritura}
\label{sec:resultados-cer}

El 19 de noviembre de 2026 se procesaron 292 imágenes con el módulo de reconocimiento de Graf-IA, que utiliza Google Cloud Vision en el modo \texttt{DOCUMENT\_TEXT\_DETECTION} con la sugerencia de idioma \texttt{es}. La prueba utilizó muestras de las particiones de prueba de tres conjuntos públicos. De OpenHand-Synth se seleccionaron 80 imágenes sintéticas con referencias en español. De UJI Pen Characters se tomaron dos muestras por cada una de 66 clases de letras y se convirtieron los trazos a imágenes. De Ticha se seleccionaron 80 renglones manuscritos en español procedentes de siete documentos históricos de Oaxaca. La selección usó la semilla 20261119 \cite{openhandSynthDataset,ujiPenCharactersV2,tichaManuscriptLines}.

El CER de cada conjunto se calculó dividiendo la suma de las distancias de edición entre el total de caracteres de referencia, según el protocolo de la Sección~\ref{sec:validacion-vision}. Antes de la comparación se normalizó el texto a Unicode NFC y se sustituyeron las secuencias de espacios y saltos de línea por un solo espacio, sin espacios al inicio ni al final. Se conservaron las mayúsculas, los acentos, la ñ y los signos de puntuación. Las respuestas sin texto se contabilizaron como eliminaciones. Una solicitud de OpenHand-Synth requirió un reintento y se incluyó una sola vez en el cálculo final.

\begin{table}[htbp]
\centering
\caption{Tasa de error por carácter del OCR por conjunto de prueba.}
\label{tab:resultados-cer}
\renewcommand{\arraystretch}{1.2}
\begin{tabular}{@{}lrrrrr@{}}
\toprule
Conjunto & Imágenes & Caracteres & Ediciones & Sin texto & CER \\
\midrule
OpenHand-Synth & 80 & 4\,684 & 644 & 0 & 13.75\,\% \\
UJI Pen Characters v2 & 132 & 132 & 82 & 32 & 62.12\,\% \\
Ticha & 80 & 3\,631 & 1\,095 & 0 & 30.16\,\% \\
\bottomrule
\end{tabular}
\end{table}

Sobre las 80 imágenes sintéticas en español, el OCR obtuvo un CER de 13.75\,\%. El filtro de calidad de OpenHand-Synth limita la variedad de casos difíciles de ese conjunto \cite{openhandSynthDataset}. En UJI, cada imagen contiene una sola letra y el OCR devolvió 32 transcripciones vacías. Las muestras de Ticha son renglones de documentos históricos. La Tabla~\ref{tab:resultados-cer} presenta una tasa por conjunto para conservar el contexto de cada medición.
